\section{Exercises}\label{sec:05-tactics:06-exercises:1}

\textbf{Key points.} A goal is hypotheses above a line, a statement to prove below it; a tactic replaces it with zero or more simpler goals. The core loop is to try the cheapest tactic, read why it failed, find structure to split or induct on, and locate the specific lemma that matches. \texttt{induction} generates one case (and hypothesis \texttt{ih}) per constructor, exactly mirroring a recursive function over the same type.

\begin{enumerate}
\def\labelenumi{\arabic{enumi}.}
\tightlist
\item
  Prove \texttt{theorem and\_\allowbreak{}comm\_\allowbreak{}tac \{P Q : Prop\} (h : P ∧ Q) : Q ∧ P := by ...} using \href{https://lean-lang.org/doc/reference/latest/Tactic-Proofs/Tactic-Reference/}{\texttt{constructor}}, \texttt{h.left}, \texttt{h.right}.
\item
  Prove \texttt{theorem nat\_\allowbreak{}mul\_\allowbreak{}zero (n : Nat) : n * 0 = 0 := by rfl}. Check whether \href{https://lean-lang.org/doc/reference/latest/Tactic-Proofs/Tactic-Reference/}{\texttt{rfl}} alone works, and if not, use \href{https://lean-lang.org/doc/reference/latest/Tactic-Proofs/Tactic-Reference/}{\texttt{induction}}.
\item
  Rewrite the \texttt{modus\_\allowbreak{}ponens} proof from Chapter 4 in tactic mode.
\item
  \texttt{rw [Nat.add\_\allowbreak{}succ]} and the rest closed the \texttt{succ} case of \texttt{my\_\allowbreak{}add\_\allowbreak{}comm}, but \texttt{rw} on its own cannot touch \texttt{a + b = b + a} before any induction is performed. Explain precisely what changes between the two situations, in terms of what is available to rewrite \emph{with}.
\item
  \texttt{nat\_\allowbreak{}mul\_\allowbreak{}zero (n : Nat) : n * 0 = 0} closes by \texttt{rfl} alone. Determine whether \texttt{theorem nat\_\allowbreak{}zero\_\allowbreak{}mul (n : Nat) : 0 * n = 0} also closes by \texttt{rfl}, and prove it by whichever means is required, justifying the choice from the recursive definition of \texttt{Nat.mul}.
\item
  A tactic that produces no error but also does not close the goal is neither a success nor a failure. Explain what it does mean, and why the goal state, not the presence or absence of a red error message, is the thing to check after every tactic.
\end{enumerate}

Solutions, \hyperref[sec:15-appendix-solutions:05-chapter-5:1]{Appendix, Chapter 5}.

With definitions, propositions, and tactics in hand, we pause for a brief rigor check before diving into groups.
